% ID: SEL-D0069-OA-E03
% ESTADO: cotejada. Segunda pasada visual y borrador compilado revisados.
% Fuente: Selectividad/sel_2009_matematicas_aplicadas/mat apli CC SS03.pdf
% SHA-256: 8cf5f487f300c573430722e40f741601da1756c5e8c13e541f697d1620373545
% Página: 1; opción A; ejercicio original 3.

\noindent\underline{Parte I}

Un turista que realiza un crucero tiene un $50\%$ de probabilidad de visitar Cádiz, un $40\%$ de visitar Sevilla y un $30\%$ de visitar ambas ciudades. Calcule la probabilidad de que:
\begin{enumerate}[label=\alph*)]
\item \textbf{(0.5 puntos)} Visite al menos una de las dos ciudades.
\item \textbf{(0.5 puntos)} Visite únicamente una de las dos ciudades.
\item \textbf{(0.5 puntos)} Visite Cádiz pero no visite Sevilla.
\item \textbf{(0.5 puntos)} Visite Sevilla, sabiendo que ha visitado Cádiz.
\end{enumerate}
\par\noindent\begin{minipage}{\linewidth}
\noindent\underline{Parte II}

El tiempo (en horas) que permanecen los coches en un determinado taller de reparación es una variable aleatoria con distribución Normal de desviación típica 4 horas.
\begin{enumerate}[label=\alph*)]
\item \textbf{(1 punto)} Se eligieron, al azar, 16 coches del taller y se comprobó que, entre todos, estuvieron 136 horas en reparación. Determine un intervalo de confianza, al $98.5\%$, para la media del tiempo que permanecen los coches en ese taller.
\item \textbf{(1 punto)} Determine el tamaño mínimo que debe tener una muestra que permita estimar la media del tiempo que permanecen en reparación los coches en ese taller con un error en la estimación no superior a una hora y media y con el mismo nivel de confianza del apartado anterior.
\end{enumerate}
\end{minipage}\par
